% ============================================================
% R01 - Rho-class · 双栏学术报告
% 文档类：rho-class（需 Overleaf 模板或自行下载）
% 编译器：pdfLaTeX（模板已内置 pdfLaTeX 优化）
% 风格：双栏，学术期刊排版，Fira Sans 标题字体
% 适合：研究报告、实验报告、技术文档、学术论文
%
% Overleaf 使用方法：
%   搜索 "rho-class" 模板，直接 Open as Template
%   或访问：https://github.com/MemoJimenez/Rho-class
%
% 原模板作者：Guillermo Jimenez (CC BY 4.0)
% 去个人品牌版本，保留完整功能结构
% ============================================================

\documentclass[9pt,a4paper,twocolumn,twoside]{rho-class/rho}
\usepackage[english]{babel}

%% 中文支持（如需中文，改用 XeLaTeX 并替换babel）
%% \usepackage[UTF8]{ctex}

%----------------------------------------------------------
% 文档元数据
%----------------------------------------------------------

\doctype{[报告类型，如：Research Article / Lab Report / Technical Report]}
\title{[报告标题]}

%----------------------------------------------------------
% 作者与机构（支持多作者）
%----------------------------------------------------------

\author[{\textasteriskcentered},a,1]{作者一}
\author[{\textdagger},b,2]{作者二}
\author[{\textdagger},c,3]{作者三}

\affil[a]{机构一}
\affil[b]{机构二}
\affil[c]{机构三}

\dates{编写于 \today}

%----------------------------------------------------------
% 通讯作者 & 文档说明
%----------------------------------------------------------

\corres{\textsuperscript{\textasteriskcentered}通讯作者：\href{mailto:author@institute.org}{author@institute.org}}
\corres{\textsuperscript{\textdagger}这些作者对本工作有同等贡献。}

\docinfo{本文档使用 Rho-class (v.3) 排版，pdfLaTeX 编译，无错误。}

%----------------------------------------------------------
% 期刊/期卷信息
%----------------------------------------------------------

\journalname{[机构/期刊名称]}
\journal{[期刊全名] 2026}
\theday{\today}
\vol{X}
\no{X}

%----------------------------------------------------------
% 摘要与关键词
%----------------------------------------------------------

\begin{abstract}
    [在此处填写摘要。描述研究背景、方法、主要结果和结论。
    建议 150--250 字。支持 \LaTeX{} 格式化命令。]
\end{abstract}

\keywords{[关键词1], [关键词2], [关键词3], [关键词4]}

%----------------------------------------------------------

\begin{document}

    \maketitle
    \thispagestyle{firststyle}

%----------------------------------------------------------
% 可选目录（默认关闭，去掉注释启用）
%----------------------------------------------------------
% \tableofcontents

%----------------------------------------------------------

\section{引言}

    \rhostart{在}此处撰写引言，用 \texttt{\textbackslash rhostart\{首字\}} 命令
    实现首字下沉效果。说明研究背景、研究问题及本文贡献。

    段落正文支持正常 \LaTeX{} 排版。引用参考文献使用
    \texttt{\textbackslash cite\{\}} 命令，如 \cite{ref1}。

%----------------------------------------------------------

\section{方法}

    \subsection{实验设计}
    
        在此描述实验方法、数据收集过程和分析手段。

    \subsection{数据分析}

        公式示例（\equref{eq:example}）：

        \begin{equation}\label{eq:example}
            i\hbar \frac{\partial}{\partial t}\Psi(\mathbf{r},t) =
            \left[\frac{-\hbar^2}{2m}\nabla^2 + V(\mathbf{r},t)\right]
            \Psi(\mathbf{r},t)
        \end{equation}

%----------------------------------------------------------

\section{结果}

    \subsection{数据表格}

        \tabref{tab:results} 展示实验数据。

        \begin{table}[H]
            \caption{实验结果数据}
            \label{tab:results}
            \begin{tabular}{
                >{\raggedright\arraybackslash}p{0.25\columnwidth}
                >{\raggedright\arraybackslash}p{0.29\columnwidth}
                >{\raggedright\arraybackslash}p{0.31\columnwidth}
            }
            \toprule
            \textbf{条件} & \textbf{测量值} & \textbf{误差} \\
            \midrule
            条件 A & 1.234 & ±0.012 \\
            条件 B & 5.678 & ±0.034 \\
            条件 C & 9.012 & ±0.056 \\
            \bottomrule
            \end{tabular}
            \tabletext{注：数值为三次重复测量的均值。}
        \end{table}

    \subsection{图示}

        \figref{fig:result} 展示实验结果图。

        \begin{figure}[H]
            \centering
            \includegraphics[width=0.9\columnwidth]{figures/example.pdf}
            \caption{实验结果图示。}
            \label{fig:result}
        \end{figure}

%----------------------------------------------------------

\section{讨论}

    对结果进行分析和解释，与相关文献对比，讨论局限性。

    使用提示框环境示例：

    \begin{note}
        这是一个说明框（note 环境）。适合放置重要提示、
        注意事项或补充说明。
    \end{note}

    \begin{rhoenv}[frametitle=自定义标题框]
        这是 rhoenv 自定义框，可以设置任意标题。
        适合放置重要定义、公式推导或代码说明。
    \end{rhoenv}

%----------------------------------------------------------

\section{代码示例}

    使用 listings 包（推荐，无需额外安装）：

    \begin{code}
        \lstinputlisting[language=python]{example.py}
        \label{lst:code1}
        \caption{Python 示例代码。}
    \end{code}

    行内代码：\inlinecode{print("Hello, World!")}

%----------------------------------------------------------

\section{结论}

    总结主要发现，说明研究的理论意义和实践价值，
    提出未来工作方向。

%----------------------------------------------------------

\section*{致谢}
\addcontentsline{toc}{section}{致谢}

    感谢 [合作者/资助方/机构] 对本研究的支持。

%----------------------------------------------------------

\printbibliography

%----------------------------------------------------------

\end{document}
